% Section 9: CONFIGURATION AND DEPLOYMENT

\section{Configuration and Deployment}

\subsection{Default Parameters}
\begin{lstlisting}[caption=System Default Configuration]
// PID Parameters (from SettingsHandler.h)
DEFAULT_SETPOINT = 3.0 bar
DEFAULT_KP = 0.0
DEFAULT_KI = 0.0  
DEFAULT_KD = 0.0

// PWM Configuration
DEFAULT_PWM_FREQ = 2000 Hz
DEFAULT_PWM_RES = 14 bits
DEFAULT_PID_SAMPLE_TIME = 10 ms

// Control System
DEFAULT_CONTROL_FREQ_HZ = 1000 Hz
DEFAULT_FILTER_STRENGTH = 0.0

// Advanced Features
DEFAULT_ANTI_WINDUP = false
DEFAULT_HYSTERESIS = false
DEFAULT_HYST_AMOUNT = 5.0%

// Valve Limits (from Constants.h)
VALVE_MIN_DUTY = 20.0%
VALVE_MAX_DUTY = 90.0%

// Sensor Calibration (from Constants.h)
SENSOR_MIN_VOLTAGE = 0.5 V
SENSOR_MAX_VOLTAGE = 4.5 V
SENSOR_MIN_BAR = 0.0
SENSOR_MAX_BAR = 10.0

// Slider Limits (from Constants.h)
SP_MIN/MAX/STEP = 0.0 / 10.0 / 0.25
KP_MIN/MAX = 0.0 / 20000.0
KI_MIN/MAX = 0.0 / 20000.0
KD_MIN/MAX = 0.0 / 500.0

// Chart Axis Defaults (from Constants.h)
Y_MIN = 0.0 bar, Y_MAX = 10.0 bar
TIME_WINDOW = 8 s
\end{lstlisting}

\subsection{Persistent Storage}
Configuration data is automatically saved to LittleFS flash storage in JSON format, ensuring settings survive power cycles and system resets.

\subsection{Build Configuration}
\begin{lstlisting}[caption=PlatformIO Configuration]
[env:arduino_nano_esp32]
platform = espressif32
board = arduino_nano_esp32
framework = arduino
lib_deps = 
    bblanchon/ArduinoJson@^7.4.1
    adafruit/Adafruit ADS1X15@^2.5.0
    https://github.com/imax9000/Arduino-PID-Library.git
    littlefs
upload_protocol = esptool
upload_speed = 921600
board_upload.before_reset = usb_reset
extra_scripts =
    pre:inject_version.py
    usb_reset_pre_upload.py
monitor_rts = 0
monitor_dtr = 0
build_flags =
    -DCORE_DEBUG_LEVEL=0
board_build.filesystem = littlefs
board_build.partitions = default_16MB.csv
\end{lstlisting}

\subsection{Testing Strategy}
\begin{itemize}
    \item \textbf{Unit Testing:} Individual class functionality
    \item \textbf{Integration Testing:} Inter-class communication
    \item \textbf{Hardware-in-Loop:} Real sensor and valve testing
    \item \textbf{Performance Testing:} Timing and memory analysis
    \item \textbf{Safety Testing:} Emergency response validation
\end{itemize}
